\section{Social facts kept in records}\label{sec:grounding}

A child learning chess protests that a piece of wood cannot jump over the pawns. The wood cannot,
and the knight can. Before the rules of chess there is no knight, only a carved horse's head on a base.
\citet{searle1995} calls rules of this kind constitutive, and sets them beside regulative rules,
which govern something that already goes on without them. Driving on the left is regulated;
checkmate exists only because rules like ``checkmate ends the game'' were written. His formula for a
constitutive rule is \emph{X counts as Y in C}: carved wood counts as a knight in chess, and printed
paper counts as a thousand rupees in Nepal. The rule assigns a status, and the status comes with
powers and duties that the physical object could never carry alone, which Searle calls status
functions. In later work he argues that one kind of act assigns every status: we make something the
case by declaring that it is the case, as in ``I now pronounce you married'' or ``the meeting is
open'', and he calls these acts status function declarations \citep[pp.~13, 93]{searle2010}.

Searle's account rests on what people accept. \citet{ferraris2013} rests his on what people write
down. A wedding that nobody records becomes, once the guests who remember it are gone, hard to tell
apart from one that never happened. For Ferraris a social object is an inscribed act: a promise
becomes a debt when it is written into an IOU, and a sale becomes a title when it is entered in a
register. He distinguishes documents that inscribe an act, such as a will or a contract, from
documents that merely register a fact, such as a temperature copied into a chart, and places social
objects with the first kind. Records let a social fact outlast the moment it was made, and carry it
to people who were never in the room.

\citet[ch.~6]{epstein2015} supplies a distinction that the rest of the paper depends on. A neighborhood
committee decides things by written proposal. At its founding meeting its members agreed on a rule:
a proposal goes to a vote once three members have endorsed it, and an endorsement is a signed record
saying so. Months later Sita files a proposal to fix the streetlight by the temple. Hari signs an
endorsement, then Maya, then Bishnu, and the proposal is ready for a vote. Epstein's distinction
separates what is at work here. Grounding is what makes a particular fact hold: these three signed endorsements make this
proposal ready for a vote. Anchoring is what put the rule in place: the founding meeting at which
the members accepted that three endorsements would be enough. Collective acceptance sets up the
frame, and records fix the facts inside it.

Under a frame of this kind, every fact has what we call a \emph{grounding condition}: a condition on
the records that says when the fact holds. \emph{The proposal has three endorsements} holds when
there are three endorsement records from three members. \emph{Hari endorsed it} holds when Hari's
signed endorsement exists. A question about authority then becomes a question about grounding
conditions: who can check a condition, and on whose records.

\section{Facts that ask for an absence}\label{sec:sort}

In an imagined clinic, a family and its clinics keep a diabetic patient's records together. The
patient's meter signs each glucose reading, the laboratory signs each result, doctors sign
prescriptions, and a family member signs when a delivery of insulin arrives. Each keeps their own
log, and records flow between them.

At two in the morning the night nurse looks at what has reached her. A reading of 7.2 was taken at
ten past ten the evening before, signed by the meter. The laboratory's result from Tuesday is in.
The last prescription she can see says ten units. Some of what she knows will hold whatever arrives
next. The reading was taken, and nothing that arrives later can untake it. The laboratory reported
its result; a correction would be a new record, and the old result would still have been reported.
Facts like these ask only that certain records exist, and adding records can never make them false.
Their grounding conditions are monotone: more input never takes away a conclusion already drawn.

One thing she knows is different. \emph{The current dose is ten units} asks for more than the
prescription she holds. It asks that there be no later prescription, and the day doctor may have
signed one at nine in the evening on a laptop that has not yet synchronized. The nurse can be careful
and competent and still give the wrong dose, because the fact she needs depends on the absence of a
record, and no amount of looking at her own records can establish an absence.

Absence facts turn up everywhere once they have a name: no objection remains open, this is the
latest version, this license has not been revoked, this slot is still free. Each holds of some set
of records and fails of a larger one, which makes its grounding condition non-monotone.\footnote{Every
grounding condition does exactly one of the two: it asks only for records or it asks for an absence.
When finitely many records lie between the smaller set and the larger, the absence comes down to a
single record that the condition needs to be missing, and over finitely many records a condition
that asks only for records is a positive Boolean expression in which records are present. We have
proved these in Lean 4, as \texttt{sort}, \texttt{absence\_step} and \texttt{positive\_iff\_posBool},
in the library \texttt{CosysSocial} of release \texttt{proofs-v1.1.0} \citep{karna2026social}, which
builds on the proofs for \citet{karna2026}. Each result below is named as it is there.}

Our claim is that the CALM theorem draws the same line through social facts that it draws through
programs, as soon as the facts are grounded in records kept in many hands. A fact with a monotone
grounding condition can be settled by anyone who holds the records, in any order, and never has to
be taken back. A fact whose grounding condition asks for an absence cannot be settled that way by
anyone, however diligent, until some coordination has happened.\footnote{In \texttt{proofs-v1.1.0}
this is \texttt{calm\_social}, stated in the network model of \citet{karna2026}: a grounding
condition asks only for records exactly when every network whose machine answers it settles the
same way under every fair order of delivery. Its conditions are given with the reply to the
objection that the theorem does not transfer (\secref{sec:obj-transfer}). The nurse's side of it is
\texttt{positive\_iff\_local\_sound}: a conclusion each observer draws from its own records holds of
all the records together exactly when its condition asks only for records.} The nurse's local fact, that she has
seen no later prescription, is monotone and safe to state. The institutional fact, that no later
prescription exists, is the one the dose depends on, and a late or partitioned record can falsify
it.
